\documentclass[12pt]{article}
\usepackage[utf8]{inputenc}
\usepackage[T1]{fontenc}
\usepackage{hyperref}
\title{Khamnigan Mongol Language Resources for K-12 Curriculum Development}
\author{Compiled by Claude Sonnet 5.5 (claude-sonnet-5-5) for Andr\'as Kornai}
\date{}
\begin{document}
\maketitle

\section{Introduction}
Khamnigan Mongol (ISO 639-3 ykh; Glottocode kham1281) is a Mongolic language spoken in China, Mongolia and the Russian Federation. It is poorly documented, and the search results describe its speakers as mostly elderly. The search found no school materials. Everything below is linguistic reference material.

\section{Primary Educational Resources}
No primers, textbooks or school programmes for Khamnigan Mongol turned up in the search. The language is not known to be taught in schools.

\section{Linguistic and Reference Materials}
\begin{itemize}
\item Khamnigan Mongol, general overview (Wikipedia): \url{https://en.wikipedia.org/wiki/Khamnigan_Mongol}
\item Glottolog languoid page for Khamnigan Mongol: \url{https://glottolog.org/resource/languoid/id/kham1281}
\item Janhunen, ``Khamnigan Mongol'' (2003), a grammar sketch in \emph{The Mongolic Languages}, listed in Glottolog: \url{https://glottolog.org/resource/reference/id/cldf:janhunen:03:1}
\item Wonsoo Yu (2011), \emph{A Study of the Mongol Khamnigan Spoken in Northeastern Mongolia}, listed in Glottolog as a grammar sketch and dictionary: \url{https://glottolog.org/resource/reference/id/39648}
\item Basic vocabularies of several Khamnigan communities in China and Mongolia: \url{https://cris.huji.ac.il/en/publications/basic-vocabularies-of-several-khamnigan-communities-in-china-mong/}
\item Basic vocabularies of three groups of Khamnigans (Altaist): \url{https://www.altaist.org/?p=5201}
\item Yamakoshi, Comparative Dictionary of Mongolic Languages, which includes Khamnigan basic vocabulary: \url{https://mongolicbv.aa-ken.jp/download/download_version_20260325.html}
\end{itemize}

\section{Religious and Community Resources}
No religious or community-organisation resources turned up in the search.

\section{Digital and Audio Resources}
No audio, video or software resources turned up in the search.

\section{Curriculum Development Context}
Documentation is limited to linguistic description and wordlists. A native-language curriculum would first need community-led documentation: recordings, an agreed orthography and teaching materials. The Khamnigan--Russian dictionary of Damdinov and Sundueva is mentioned in the search summary, but no URL for it was seen.
\end{document}